\documentclass[12pt]{article}
\usepackage[utf8]{inputenc}
\usepackage[T1]{fontenc}
\usepackage{hyperref}
\title{Kuku Language Resources for K-12 Curriculum Development}
\author{Compiled by Claude Sonnet 5 (claude-sonnet-5) for Andr\'as Kornai}
\date{}
\begin{document}
\maketitle

\section{Introduction}
Kuku (ISO 639-3 ukv; Glottocode kuku1285) is an Eastern Nilotic language closely related to Bari and Kakwa. Search results describe it as spoken mainly in Kajo Keji County, Central Equatoria, South Sudan, with a community in the West Nile region of northern Uganda. Speaker figures vary across sources. Some sources treat it as a dialect of Bari, and it has no standardized orthography; it is sometimes written with Bari or Kakwa conventions. Kakwa (Yei, Koboko) is a different language and its resources should not be mixed in.

\section{Primary Educational Resources}
No Kuku primers, textbooks or school programmes were found in search results.

\section{Linguistic and Reference Materials}
\begin{itemize}
\item Wikipedia, ``Kuku dialect'': \url{https://en.wikipedia.org/wiki/Kuku_dialect}
\item Wikipedia, ``Kuku people'': \url{https://Www.Wikipedia.org/wiki/Kuku_people}
\item Omniglot, Kuku: \url{https://www.omniglot.com/writing/kuku.htm}
\end{itemize}

\section{Religious and Community Resources}
\begin{itemize}
\item Joshua Project language page for ukv: \url{https://joshuaproject.net/languages/ukv}
\end{itemize}

\section{Digital and Audio Resources}
None found.

\section{Curriculum Development Context}
The language is not known to be taught in schools. Orthography is unstandardized, and the closely related Bari has more written tradition that might be adapted, but this was not examined here. Community consultation in Kajo Keji and the West Nile region would be required.

\end{document}